\section{Introduction}
\label{sec:introduction}
CLIP enables image classification by comparing images with natural-language descriptions of classes, without task-specific training labels~\citep{radford2021clip}. When a small labelled support set becomes available in a target domain, a linear classifier or a cache-based adapter can exploit it. However, higher average accuracy does not reveal whether adaptation also breaks predictions that were already correct. Such correct-to-incorrect changes are \emph{negative flips}~\citep{yan2021flips}. For a user deciding whether to replace a zero-shot classifier, both the corrected errors and the newly introduced errors matter.

We study this trade-off in few-shot domain adaptation on Office-Home~\citep{venkateswara2017officehome}. The image and text encoders of CLIP ViT-B/32 remain frozen. Each of the four domains serves as the target, with $K\in\{1,2,4,8\}$ labelled images per class and five support-set seeds. We compare an ordinary linear probe (A1), a probe initialized from zero-shot text prototypes (A2), and training-free Tip-Adapter (A4)~\citep{zhang2022tip}. Hyperparameters are selected by simulating the task on the other three domains; target test labels are reserved for evaluation. Thus, the setting includes target support labels but no target validation set.

The study addresses three questions: \textbf{RQ1}, how much of adaptation's net accuracy change comes from correcting errors versus introducing negative flips; \textbf{RQ2}, how these changes vary with zero-shot confidence; and \textbf{RQ3}, whether choosing between zero-shot and adapted predictions per image improves the accuracy--harm trade-off. We evaluate a confidence gate and a few-shot, frozen-feature variant of VRF-style distance weighting~\citep{zhu2024vrf}, with fixed probability mixing and matched random controls. Sample-wise combination already appears in VRF and entropy-based weight interpolation in DaWin~\citep{oh2025dawin}; our contribution is its controlled evaluation in this limited-label setting, together with a flip-level diagnosis.

The results favour retaining zero-shot information, but offer little support for the additional complexity of sample-wise selection. At $K=4$, A2 improves mean accuracy from 79.63\% to 82.99\%, while harming 3.74\% of test images. Fixed probability mixing raises accuracy to 83.24\% and lowers harm to 2.43\%. Confidence gating beats matched random selection, yet adds little over A2 alone; VRF-style weighting is also close to fixed mixing and to its random-set control. The analysis explains this pattern: corrections and negative flips overlap in the low-confidence part of the test set, while the few-shot failure-set distance is a weaker error signal than zero-shot confidence.
